\documentclass[12pt,letterpaper]{article}
\usepackage[spanish]{babel}
\usepackage{fontspec}
\setmainfont{Arial}
\usepackage[margin=2.54cm]{geometry}
\usepackage{xcolor,graphicx,hyperref,setspace,tikz}
\definecolor{industrialBlue}{HTML}{123B5D}
\definecolor{industrialOrange}{HTML}{F05A28}
\hypersetup{colorlinks=true,linkcolor=industrialBlue,urlcolor=industrialBlue}
\onehalfspacing
\begin{document}
\begin{center}
{\large\bfseries Instituto Tecnologico Superior de Cajeme\par}
{\large Maestria en Gestion Administrativa\par}
\vspace{0.8cm}
{\LARGE\bfseries\color{industrialBlue} Imagen de la empresa\par}
\vspace{0.25cm}
{\Large Industrial Revolucionaria\par}
\vspace{0.7cm}
\begin{tabular}{rl}
Asignatura: & Plan de Negocios (GGPN01) \\
Modulo: & 6545 \\
Estudiante: & Martin Jonathan de la Cruz Muñoz \\
Periodo: & 2026-II \\
Fecha: & 21 de septiembre de 2026 \\
\end{tabular}
\end{center}

\vspace{0.8cm}
\begin{center}
\begin{tikzpicture}[scale=0.95]
  \foreach \angle in {0,40,...,320} {\draw[industrialBlue,line width=5pt] (0,0) -- (\angle:1.2);}
  \fill[industrialBlue] (0,0) circle (0.95);
  \fill[white] (0,0) circle (0.72);
  \node[text=industrialBlue,font=\bfseries\Large] at (0,0) {IR};
  \node[anchor=west,text=industrialBlue,font=\bfseries\Large] at (1.5,0.4) {INDUSTRIAL};
  \node[anchor=west,text=industrialOrange,font=\bfseries\Large] at (1.5,-0.15) {REVOLUCIONARIA};
  \draw[industrialOrange,line width=1.5pt] (1.5,-0.42) -- (6.6,-0.42);
  \node[anchor=west,text=black!70,font=\small] at (1.5,-0.75) {Mantenimiento que impulsa tu operacion};
\end{tikzpicture}
\end{center}

\section*{Nombre del proyecto}
El nombre \textbf{Industrial Revolucionaria} identifica una empresa de mantenimiento y soluciones industriales para micro, pequenas y medianas empresas de Monterrey y su zona metropolitana. La palabra \textit{Industrial} comunica el sector atendido; \textit{Revolucionaria} expresa la intencion de transformar el mantenimiento reactivo en un servicio planificado, documentado y orientado a la mejora continua.

El nombre se eligio en espanol para facilitar su pronunciacion y recordacion en el mercado local. Es congruente con el proposito del proyecto: reducir paros operativos mediante diagnostico, mantenimiento preventivo y correctivo, seguimiento de ordenes de trabajo y gestion de refacciones. Su disponibilidad y recordacion deben validarse antes de cualquier uso comercial.

\section*{Logotipo propuesto}
El logotipo fue diseñado por el estudiante como propuesta original para esta actividad. El engrane simplificado representa maquinaria, precision tecnica y continuidad operativa; el monograma \textbf{IR} permite identificar de manera compacta el nombre en ordenes de servicio, uniformes, vehiculos y medios digitales.

El azul industrial (\texttt{\#123B5D}) se eligio para comunicar confianza, orden y profesionalismo. El naranja de seguridad (\texttt{\#F05A28}) funciona como acento que representa energia, atencion y respuesta ante contingencias. La combinacion mantiene contraste y legibilidad; su aceptacion con clientes potenciales debe comprobarse posteriormente.

\section*{Lema y utilidad estrategica}
El lema propuesto es \textbf{``Mantenimiento que impulsa tu operacion''}. Surge de entender el mantenimiento como una actividad que contribuye a la continuidad, seguridad y eficiencia de las empresas clientes, no solo como reparacion de fallas. El verbo ``impulsa'' vincula el servicio con un resultado operativo sin prometer eliminar por completo las contingencias.

Estrategicamente, el lema sintetiza la propuesta de valor y complementa el logotipo en presentaciones, cotizaciones, vehiculos y medios digitales. Tambien diferencia al proyecto de servicios comunicados unicamente como respuesta de emergencia. Su comprension y capacidad de diferenciacion deben validarse mediante entrevistas con responsables de mantenimiento, planta y compras.

\section*{Validacion pendiente}
La identidad es una propuesta academica. Antes de adoptarla comercialmente se verificara la disponibilidad del nombre, la originalidad del signo, su legibilidad en escala de grises, su reproduccion en distintos materiales y la opinion de clientes potenciales. El formato solicitado por el modulo es Word; el PDF es una copia de revision y no lo sustituye.
\end{document}